% Transcription — fonds Alexandre Grothendieck, shelfmark 133, batch 1 (pages 1-20).
% Established from the facsimile archives/batches/133/batch-01.pdf (a mirror of
% https://grothendieck.umontpellier.fr/133.pdf), per the skill
% transcribe-grothendieck. The batch file carries no cover sheet: PDF page N
% is page N of the folder (the Montpellier footer prints « 1 » ... « 20 »),
% and the archivists' pencil numbers bottom-left agree wherever legible
% (« 3 », « 4 », « 5 » ... « 20 »).
%
% Pass: Opus 5.5 (claude-opus-5-5), 2026-09-23 — first pass, unchecked against the pages by a human.
%
% Ten of the twenty pages are transcribed. Absent: page 1, an otherwise
% blank sheet of pink card inscribed in his hand « groupes Witt et formes
% quadratiques » — a cover, whose words become the first \section heading
% (references/hand.md §5); pages 4 and 9, blank; pages 10-16, two offprints,
% printed, with nothing in his hand on any of the seven leaves (checked at
% 250 dpi and by an ink-colour count), hence absent under folder 22's rule:
%   - O. Laborde, « Théorème de Skolem-Nother pour les algèbres d'Azumaya
%     graduées sur un anneau semi-local », C. R. Acad. Sc. Paris, t. 279
%     (16 septembre 1974), Série A, p. 447-449 (pages 10-12);
%   - O. Laborde, « Formes quadratiques, algèbres de Clifford et
%     signatures », C. R. Acad. Sc. Paris, t. 278 (24 juin 1974), Série A,
%     p. 1599-1602 (pages 13-16).
% They are identified once, in a \note before page 17, because page 7 cites
% « Laborde dixit ». He did not paginate these leaves.
%
% What the batch is.
%   2      Scratch on a sheet of USTL Institut de Mathématiques letterhead:
%          the associativity square for a tensor product of complexes L.,
%          an obstruction in Ext^1(π_0,π_0,π_0; π_1) hit by δ^(1) from
%          Ext(π_0,π_0; π_1), and Mac Lane's pentagon for x, y, z, t.
%   3      Scratch: P' = W*(L^{⊗2}), η ∈ H^1(P', F_2) with
%          η^2 + w_1 η + w_2 = 0, w_1 = α = ∂(-1), w_2 ≡ c_1 mod 2; a
%          spectral sequence Ext_B^p(Tor_q^A(M,B), N). A few non-mathematical
%          lines at the foot of the sheet are not transcribed.
%   5-8    « Groupes de Witt et variantes » (his title): for a ringed topos
%          S, WG(S) the Grothendieck group of strictly non-degenerate
%          quadratic modules; a pre-λ-ring augmented to H^0(S,Z); the
%          reduction of the λ-ring identities to WG of the classifying topoi
%          B_{O(m)×O(n)}, B_{O(m)}; the hyperbolic map α: K_s(S) → WG(S),
%          E ↦ E ⊕ Ě, and W(S) = coker α; (M,q) ⊕ (M,-q) hyperbolic when 2
%          is invertible, and « Laborde dixit » without that condition on an
%          affine scheme, but not on B_{O(n)}; W(S) a ring augmented only to
%          H^0(S, Z/2Z); then, O_S strictly local with 2 invertible, the
%          total Stiefel-Whitney class w: WG(S) → H*(S, F_2) and a
%          λ-homomorphism WG(S) → H^0(S,Z) × (1 + Ĥ*+(S, F_2)). The attempt
%          to factor w through W(S) is struck out.
%   17-20  Black felt pen, a separate run: k_n(G) = k(G × S_n), the graded
%          k_0(G)-algebra k_*(G) of virtual finite G-sets with the induction
%          product, τ^n(E) = cl_n(Mon(I_n, E)), the exponential property
%          τ*(E ⊔ F) = τ*(E) τ*(F) with its proof, the corollary extending
%          τ* to k_0(G), and the maps q_H, q^n_H for H ⊂ S_n. Page 18 is the
%          scratch sheet beside it: λ_t, σ_t, λ^i(E) = cl(P_i(E)),
%          σ^i(E) = cl(Sym^i(E)) for a finite G-set E, and 2λ^2(E)+E = E^2.
%          Page 17 is a sheet of computer listing (printed show-through,
%          not transcribed) carrying only the formula K_n(G) = K(G × S_n).
%
% Notation fixed for the batch. L_\bullet for his « L. » on page 2; WG, W,
% K_s, K as he writes them; \check{E} for the dual; \alpha for the
% hyperbolic map (he also writes φ once, page 6, kept); B_{O(m)} for the
% classifying topos; \mathfrak{S}_n for his gothic S; \mathfrak{P}_i for a
% closed-topped F-like glyph that page 20 makes unambiguous (subsets of
% cardinality i, J ∈ P_i(I_n)) and that recurs on page 18; \mathrm{cl}_n for
% his « cl_n »; \mathrm{Mon} for « Mon ». His « loc. », « t.f. », « str. »,
% « ssi », « s-groupe » are kept. Square brackets in the text are his.
%
% The dating is the inventory's own for the folder, copied verbatim; the
% title in \foldertitle is the archivists' (square brackets), not his.
\documentclass[11pt,a4paper]{article}
% Le préambule commun à toutes les transcriptions du fonds Grothendieck.
%
% Il tient en une idée : **l'incertitude se compose, elle ne se résout pas.**
% Une transcription qui lisse un mot douteux a détruit la seule chose pour
% laquelle on la faisait — un lecteur doit pouvoir savoir, à chaque endroit,
% ce qui était sur le feuillet et ce qui a été ajouté. Les six macros qui
% suivent sont donc l'appareil critique, et non de la décoration.
%
% Les mêmes commandes sont reconnues par `scripts/render.mjs`, qui produit la
% vue de lecture du site. Une macro ajoutée ici doit l'être aussi là-bas,
% faute de quoi le rendu échouera bruyamment — ce qui est voulu : un rendu
% silencieusement faux est pire qu'un rendu absent.

\ProvidesPackage{grothendieck}[2026/08/08 Transcriptions du fonds Grothendieck]

\usepackage{amsmath,amssymb,amsthm}
\usepackage{mathrsfs}
\usepackage{stmaryrd}
% Les diagrammes commutatifs sont partout dans ce fonds : c'est la langue
% même de la géométrie algébrique de Grothendieck, et une transcription qui
% les rendrait en prose ne transcrirait rien.
\usepackage{tikz-cd}
% Français par défaut — c'est la langue des feuillets — mais l'anglais est
% chargé aussi : la traduction et le résumé sont des documents anglais, et la
% typographie française (espaces avant les deux-points) y serait une faute.
% Ces documents basculent par \selectlanguage{english} en tête de corps.
\usepackage[english,french]{babel}
\usepackage{fontspec}
\usepackage{geometry}
\geometry{a4paper,margin=25mm,marginparwidth=28mm}
\usepackage{marginnote}
\usepackage{xcolor}
% ulem, not soul. `soul`'s \st and \ul are built on a character-by-character
% scan that XeLaTeX's font machinery breaks: the document fails with an
% undefined control sequence rather than degrading. ulem does the same two jobs
% and is Unicode-engine safe. `normalem` stops it hijacking \emph, which the
% transcriptions use for Grothendieck's own emphasis.
\usepackage[normalem]{ulem}
\usepackage{hyperref}
\hypersetup{colorlinks=true,linkcolor=black,urlcolor=black,pdfborder={0 0 0}}

% --- Métadonnées du lot -----------------------------------------------------
% Reprises telles quelles par le script de rendu, qui les affiche en tête de
% la vue de lecture. Elles fixent ce que le fichier prétend couvrir : sans
% elles, une transcription flotte sans point d'attache dans un fonds de seize
% mille feuillets.

\newcommand{\folder}[1]{\gdef\grothFolder{#1}}
\newcommand{\batch}[1]{\gdef\grothBatch{#1}}
\newcommand{\leaves}[2]{\gdef\grothFirst{#1}\gdef\grothLast{#2}}
\newcommand{\foldertitle}[1]{\gdef\grothFoldertitle{#1}}
\gdef\grothFolder{?}\gdef\grothBatch{?}\gdef\grothFirst{?}\gdef\grothLast{?}\gdef\grothFoldertitle{}

% --- Le résumé --------------------------------------------------------------
%
% Il ouvre la lecture modernisée, sous le titre « Résumé », et il est composé
% en petit. Ce n'est pas un ornement : c'est le seul passage du document qui
% ne soit pas des mathématiques, et un lecteur doit pouvoir le reconnaître
% avant de l'avoir lu — puis le sauter s'il connaît déjà le sujet, ou s'y
% arrêter s'il ne le connaît pas. Le filet qui le suit marque la reprise du
% corps.

\newenvironment{resume}
  {\small\setlength{\parskip}{0.45em}}
  {\par\medskip\hrule\medskip}

% --- Mots-clés ---------------------------------------------------------------
%
% En anglais, en clôture du résumé de la lecture modernisée : le vocabulaire
% moderne sous lequel chercher ce que les feuillets construisent. Le manifeste
% du site (`npm run manifest`) les extrait pour étiqueter la cote entière —
% cette ligne est la source unique des tags, il n'y a pas de fichier à côté.
\newcommand{\keywords}[1]{\par\smallskip\noindent
  {\footnotesize\textsc{Keywords} --- \textit{#1}}\par}

% --- L'appareil critique ----------------------------------------------------

% Le numéro de feuillet, en marge. C'est l'ancre qui permet de régler un
% désaccord sur une formule en regardant *un* feuillet plutôt qu'une chemise
% de six cents. Le site s'en sert aussi pour tourner les pages du fac-similé
% quand on fait défiler la transcription.
\newcommand{\leaf}[1]{%
  \marginnote{\footnotesize\color{gray!70}\textsf{#1}}%
  \label{leaf:#1}\ignorespaces}

% Illisible : on ne devine pas. Un mot inventé qui se lit comme les autres est
% le pire résultat possible.
\newcommand{\ill}{\textcolor{red!60!black}{[\,\ldots\,]}}

% Lecture douteuse : le mot est donné, et signalé comme incertain.
\newcommand{\uncertain}[1]{\uline{#1}}

% Ajout de l'éditeur — une lettre manquante, un mot sous-entendu.
\newcommand{\add}[1]{\textcolor{blue!50!black}{[#1]}}

% Biffé par Grothendieck lui-même. Ce qu'il a rayé fait partie du document :
% une rature dit souvent où la pensée a changé de direction.
\newcommand{\struck}[1]{\sout{#1}}

% Note du transcripteur, sur l'état matériel du feuillet.
\newcommand{\note}[1]{\par\smallskip{\small\color{gray!60!black}\textit{[#1]}}\par\smallskip}

% Note marginale de Grothendieck, distincte du corps du texte.
\newcommand{\marginal}[1]{\par\smallskip\noindent\hspace{1em}%
  {\small\color{gray!60!black}#1}\par\smallskip}

% --- Titre ------------------------------------------------------------------

\AtBeginDocument{%
  \begin{center}
    {\large\bfseries Fonds Alexandre Grothendieck --- cote n\textsuperscript{o}~\grothFolder}\\[2pt]
    {\itshape\grothFoldertitle}\\[4pt]
    {\small feuillets \grothFirst--\grothLast\ (lot \grothBatch)}\\[2pt]
    {\footnotesize\color{gray!60!black}Transcription établie d'après le fac-similé numérisé
     par l'Université de Montpellier.\\
     Les lectures incertaines sont soulignées, les illisibles notées [\,\ldots\,],
     les ajouts entre crochets.}
  \end{center}
  \vspace{1em}}

\endinput


\folder{133}
\batch{1}
\pages{1}{20}
\dating{1974-[à partir de 1975]}
\watermark{Édition de démonstration}
\foldertitle{[Groupes de Witt et formes quadratiques, groupes formels] : notes manuscrites (s.d.), tirés à part (1974).}

\begin{document}

\section{Groupes Witt et formes quadratiques}
\note{titre de sa main, en haut à droite du feuillet 1, une chemise de
carton rose par ailleurs vierge qui ne reçoit pas de numéro de page ; le
dernier mot est abrégé, « qu.dr.tiques », et lu \uncertain{quadratiques}}

\page{2}
\note{brouillon, sans phrase ; en tête du feuillet, deux relations isolées :}
\[
  q' + q'' \;\uncertain{+}\; q''' = -1 \qquad\qquad q''' = q'q''\,\ill{}
\]
\note{le signe entre $q''$ et $q'''$ est surchargé ; après $q'q''$, une
lettre ou un indice illisible}

\struck{\ill{}}
\note{deux lignes biffées de grands traits obliques, dont la seconde se
termine par « $(x \otimes y) \otimes z \simeq x \otimes (y \otimes z)$ »,
lisible sous la rature}

\struck{$L_{\bullet} \otimes L_{\bullet} \otimes L_{\bullet}$}

\begin{tikzcd}
  (L_{\bullet} \otimes L_{\bullet}) \otimes L_{\bullet} \arrow[rr, "a"] \arrow[d, "p \otimes L_{\bullet}"'] & & L_{\bullet} \otimes (L_{\bullet} \otimes L_{\bullet}) \arrow[d, "L_{\bullet} \otimes p"] \\
  L_{\bullet} \otimes L_{\bullet} \arrow[dr, "p"'] & & L_{\bullet} \otimes L_{\bullet} \arrow[dl, "p"] \\
  & L &
\end{tikzcd}

\note{entre les deux $L_{\bullet} \otimes L_{\bullet}$, une flèche courbe
marquée \uncertain{$\alpha$}, biffée}

\[
  \mathrm{Ext}^{1}(\pi_0, \pi_0, \pi_0 ; \pi_1)
  \xleftarrow{\ \delta^{(1)}\ }
  \mathrm{Ext}(\pi_0, \pi_0 ; \pi_1)
\]
\note{après $\pi_1$, dans le premier $\mathrm{Ext}$, une suite biffée où se
lit « $L_{\bullet} \otimes L_{\bullet}$ » ; l'exposant du second
$\mathrm{Ext}$, s'il y en a un, ne se lit pas}

\begin{tikzcd}[column sep=small]
  ((x \otimes y) \otimes z) \otimes t \arrow[r] \arrow[d] & (x \otimes y) \otimes (z \otimes t) \arrow[dd, leftrightarrow] \\
  (x \otimes (y \otimes z)) \otimes t \arrow[d] & \\
  x \otimes ((y \otimes z) \otimes t) \arrow[r] & x \otimes (y \otimes (z \otimes t))
\end{tikzcd}

\page{3}
\[
  P' = W^{*}(L^{\otimes 2}) \qquad
  L_{P'}^{\otimes 2} \simeq \mathcal{O}_{P'}
\]
\note{la lettre $W$ porte un astérisque en exposant, au-dessus de la
parenthèse}

\[
  \eta \in H^{1}(P', \mathbb{F}_2)
\]
\note{écrit sous une accolade qui embrasse les deux formules précédentes}

\[
  \eta^{2} + w_1 \eta + w_2 = 0
\]

Je dis que
\[
  w_1 = \alpha \overset{\mathrm{df}}{=} \partial(-1), \qquad
  w_2 \equiv c_1 \bmod 2
\]
\[
  \partial : H^{0}(S, \mathbb{G}_m) \to H^{1}(S, \mu_2)
\]

\[
  E = L + L^{-1} \qquad
  Q(x, x') = \struck{\ill{}}\ x x'
\]

\[
  \prod_i (1 + \alpha + \xi_i)
\]

\[
  (1 + y)(1 + y + \alpha) = 1 + \alpha + \underbrace{y(y + \alpha)}_{\overset{\shortparallel}{c_1\,?}}
\]

\[
  A \to B, \quad M, \ N \qquad
  E_2^{p,q} = \mathrm{Ext}_B^{p}(\mathrm{Tor}_q^{A}(M, B), N)
\]
\note{les indices $p$ et $q$ sont écrits par-dessus un $B$ et un $A$
d'abord posés à leur place ; à côté, un petit croquis de quadrant avec
quelques points}

\subsection{Groupes de Witt et variantes}
\note{titre de sa main, en tête du feuillet 5}

\page{5}
$S$ topos loc.\ annelé

$WG(S) =$ groupe enveloppant du monoïde commutatif des classes d'isomorphie
des modules quadratiques sur $S$ [loc.\ libres de t.f.] \emph{strictement}
non dégénérés (i.e.\ la forme bil.\ ass.\ est non dégénérée)

\emph{NB} $WG(S)$ est un pré-$\lambda$-anneau (sans unité en général…)
augmenté vers “l'anneau binomial” $H^{0}(S, \mathbb{Z})$ ($\simeq \mathbb{Z}$
si $S$ connexe). Si $S$ connexe \struck{\ill{}} et $2$ non inversible dans
$\mathcal{O}_S$, alors $WG(S) \xrightarrow{\varepsilon} \mathbb{Z}$ prend ses
valeurs dans $2\mathbb{Z}$ (les modules quadratiques str.\ non dég.\ sont de
rang pair). Si $2$ est inversible dans $\mathcal{O}_S$, alors $WG(S)$ est un
anneau avec unité. Même dans ce cas j'ignore si c'est un $\lambda$-anneau
(et pas seulement un pré-$\lambda$-anneau), donc si on a les formules
universelles
\[
  \lambda^{i}(xy) = P_i(x, \struck{\ill{}}, \ldots, \lambda^{i-1}x ; y, \struck{\ill{}}, \ldots, \lambda^{i-1}y) \qquad (i \geqslant 2)
\]
\[
  \lambda^{i}(\lambda^{j} x) = Q_{i,j}(x, \lambda^{2}x, \ldots, \lambda^{ij}(x)) \qquad (i, j \geqslant 2)
\]
($P_i$, $Q_{ij}$ pol.\ univ.\ à coeff.\ dans $\mathbb{Z}$…)
\note{dans la première formule, le deuxième argument de chaque liste est
noyé sous l'encre ; le dernier se lit $\lambda^{i-1}$, non $\lambda^{i}$, ce
qui est ainsi transcrit}

\emph{NB} Pour la première formule p.\ ex., il suffit de se placer sur le
topos classifiant d'un groupe $O(m) \times O(n)$, en prenant pour $x$ et $y$
les classes des images inverses sur $B_{O(m) \times O(n)}$ des
\struck{\ill{}} modules quadratiques universels \struck{\ill{}} de rang $m$
resp.\ $n$ sur $B_{O(m)}$ resp.\ $B_{O(n)}$. Cela nous amène donc à étudier
$WG(B_{O(m) \times O(n)})$. De même la deuxième formule à
\struck{\uncertain{vérifier}} \uncertain{tester} dans un $B_{O(m)}$, d'où la
tâche de déterminer $WG(B_{O(m)})$.

\page{6}
On a un hom.\ de groupes
\[
  K_s(S) \xrightarrow{\ \alpha\ } WG(S)
\]
en associant à la classe des \struck{\ill{}} Modules loc.\ libres de t.f.\ $E$
le \struck{\ill{}} module quadratique $E \oplus \check{E}$, avec la forme
$q(x, x') = \langle x, x' \rangle$ ; ici $K_s(S)$ désigne le groupe $K$
“spécial”, groupe enveloppant du monoïde des classes d'isomorphie de modules
loc.\ lib.\ de t.f.\ sur $S$ ; il n'est pas clair que l'hom.\ précédent se
factorise par $K(S)$ (déduit de $K_s(S)$ en tenant compte des suites exactes
courtes…)

(Donc si $E$ est ext.\ non splittée de $E'$ par $E''$, il n'est pas clair si
on a $\varphi(E) = \varphi(E') + \varphi(E'')$ dans $WG(S)$ ….)
\note{« non splittée » est ajouté au-dessus de la ligne ; il écrit ici
$\varphi$ pour l'homomorphisme qu'il vient d'appeler $\alpha$}

Par définition, $W(S)$ (groupe de Witt) est le conoyau de l'hom.\ précédent
\[
  K_s(S) \xrightarrow{\ \alpha\ } WG(S) \longrightarrow W(S) \longrightarrow 0
\]

\emph{NB} Si $2$ inv.\ dans $\mathcal{O}_S$, alors pour tout \struck{\ill{}}
module quadratique $(M, q)$, la somme $F = (M, q) \oplus (M, -q)$ définit
\struck{\ill{}} l'élément nul de $W(S)$, i.e.\ est $\simeq$ à un
$E \oplus \check{E}$ comme dessus : il suffit de prendre $E = M$, immergé
dans $F$ diagonalement et antidiagonalement [en fait, les mod.\ quad.\ $F$ de
la forme $(M, q) \oplus (M, -q)$ sont, si on veut, \struck{\ill{}} de la
forme $E \oplus \check{E}$…, mais \uncertain{où} $E$ est un Mod.\ loc.\ lib.\
de t.f.\ qui admet (au moins une) forme quadratique (str.)\ non
\emph{dégénérée}.
\note{le crochet ouvert ici n'est pas refermé sur la page}

De même, si $S$ est un schéma affine, il

\page{7}
paraît (Laborde dixit) que (sans condition que $2$ soit inversible) les
modules $(M, q) \oplus (M, -q)$ sont de la forme $E \oplus \check{E}$.
\note{au-dessus de « les », en petit, « $\alpha(E) =$ », ajout interlinéaire
qui se rapporte vraisemblablement à « $E \oplus \check{E}$ » : de la forme
$\alpha(E) = E \oplus \check{E}$}
(Mais ce n'est pas vrai dans \struck{\ill{}} les cas d'un topos loc.\ annelé,
p.\ ex.\ ce n'est pas vrai pour le topos classifiant $B_{O(n)}$, si
$n \geqslant 1$)
\note{au-dessus de la ligne, un « Donc » biffé}

\marginal{à vérifier}
\note{en marge gauche, avec un trait vertical qui tient tout le paragraphe
suivant}

\emph{Dans} l'un et l'autre cas, $W(S)$ peut dès lors se construire
directement comme quotient de l'ens.\ des classes d'isom.\ de modules
quadratiques, avec $(M, q) \sim (M', q')$ ssi \struck{$(M, q) \oplus (M',
-q')$ est isom.\ à un $E \oplus \check{E} = \alpha(E)$} $\exists$ mod.\ loc.\
lib.\ de t.f.\ $E$, $F$ avec
\[
  (M, q) \oplus (M', -q') + \alpha(F) \simeq \alpha(E)
\]

[\emph{NB} Si $F$ est isomorphe à un facteur direct de $E$, et si on dispose
d'un “th.\ de Witt” — ce qui est \uncertain{toujours} le cas si $S$ est le
spectre d'un \struck{\ill{}} \uncertain{corps} — alors cette condition
\emph{signifie} aussi que $\exists\, G$ loc.\ lib.\ de t.f.\ tel que
\[
  (M, q) \oplus (M', -q') \simeq \alpha(G) \ ]
\]

On voit que $\alpha(K(S))$ est un \emph{idéal} dans $WG(S)$ : donc $W(S)$ est
encore un anneau. Mais il n'est plus \emph{augmenté} vers $H^{0}(S,
\mathbb{Z})$, seulement vers $H^{0}(S, \mathbb{Z}/2\mathbb{Z})$
($\mathbb{Z}/2\mathbb{Z}$ souscrit $\mathbb{F}_2$) (donc vers $\mathbb{F}_2$
si $S$ connexe $\neq \emptyset$). Mais que peut-on dire de la
pré-$\lambda$-structure de $WG(S)$ — passerait-elle au quotient ?
\note{il écrit ici $\alpha(K(S))$, non $\alpha(K_s(S))$}

\page{8}
On suppose maintenant pour simplifier que l'Anneau $\mathcal{O}_S$ est
\emph{strict.} \emph{local} (p.\ ex.\ anneau \uncertain{structural} d'un
topos \uncertain{ét.}). On suppose de plus $2$ inv.\ dans $\mathcal{O}_S$.
On définit alors
\note{la parenthèse et la seconde phrase sont des ajouts, la seconde dans
la marge gauche, renvoyée à sa place par un signe}
\[
  WG(S) \xrightarrow{\ w\ } H^{*}(S, \mathbb{F}_2)
\]
(classe de Stiefel-Whitney), caractérisée (lorsque $S$ variable) par
\begin{enumerate}
  \item[a)] \emph{fonctorialité} en $S$
  \item[b)] $w(M, q)$ dans le cas où $M$ de rang $1$ (normalisation)
  \item[c)] $w(x + y) = w(x)\, w(y)$, $w(x) \in 1 + H^{+}(S, \mathbb{F}_2)$
  i.e.\ $w^{0}(x) = 1$.
\end{enumerate}

\struck{Un calcul facile montre que $w \circ \alpha = 1$, donc $w$ se
factorise en}
\struck{$W(S) \xrightarrow{\ w\ } 1 + \hat{H}^{*+}(S, \mathbb{F}_2)$}
\note{passage barré de grands zigzags ; dans la marge gauche, biffée avec
lui, une note oblique dont se lisent seulement « $w(\alpha(E)) =$ » et un
« $\mathcal{O}^{2}$ »}

\struck{\emph{NB} $w(M, q)$ \ill{}}

\emph{NB} Les arguments habituels de splitting montrent que $w(M \otimes N)$
et $w(\lambda^{i} M)$ se calculent par les formules
\struck{habituelles} universelles en termes de $w(M)$, $w(N)$ et $r(M)$,
$r(N)$ (rangs des Modules $M$, $N$). Donc on a un $\lambda$-hom.
\[
  WG(S) \longrightarrow H^{0}(S, \mathbb{Z}) \times (1 + \hat{H}^{*+}(S, \mathbb{F}_2))
\]
\note{devant la parenthèse, un $H$ biffé ; « $\lambda$- » est ajouté
au-dessus de « un »}

\section{[Anneau des $G$-ensembles finis virtuels et opérations $\lambda$]}
\note{titre entre crochets, de l'éditeur : les feuillets 17 à 20, au feutre
noir, ne portent pas de titre. Entre le feuillet 8 et celui-ci, les
feuillets 10 à 16 sont les tirés à part de deux notes d'O. Laborde aux
Comptes rendus, non annotés et non transcrits : « Théorème de
Skolem-Nother pour les algèbres d'Azumaya graduées sur un anneau
semi-local », C. R. Acad. Sc. Paris, t. 279 (16 septembre 1974), Série A,
p. 447-449 ; « Formes quadratiques, algèbres de Clifford et signatures »,
C. R. Acad. Sc. Paris, t. 278 (24 juin 1974), Série A, p. 1599-1602.}

\page{17}
\[
  K_n(G) = K(G \times \mathfrak{S}_n)
\]
\note{seule écriture de sa main sur le feuillet, en haut à droite}

\page{18}
\[
  \lambda_t(E) = \sum \lambda^{i}(E)\, t^{i} \qquad
  \lambda^{i}(E) = \mathrm{cl}(\mathfrak{P}_i(E))
\]
\[
  \sigma_t(E) = \sum \sigma^{i}(E)\, t^{i} \qquad
  \sigma^{i}(E) = \mathrm{cl}(\mathrm{Sym}^{i}(E))
\]
\note{le signe transcrit $\mathfrak{P}_i$ est un F à barre fermée ; la
page 20 l'emploie pour l'ensemble des parties à $i$ éléments,
$J \in \mathfrak{P}_i(I_n)$}

\[
  \struck{\sigma_t(E) = \lambda_{\frac{t}{1-t}}(E)} \qquad
  \boxed{\sigma_{-t}(E)\, \lambda_t(E) = 1}
\]
\note{la formule biffée est lue sous la rature, sans certitude}

\[
  \sigma^{i}(E) = \struck{\ill{}}\ Q_i(\lambda^{j}(E), 1 \leqslant j \leqslant i)
\]

\note{au milieu du feuillet, isolés : « $G$ \quad $E$ » et, à côté,
$E$ au-dessus de $X$ relié par un trait vertical, avec $G$ devant}

\[
  \sigma_t(L) = \sum L^{i} t^{i} = (1 - tL)^{-1} = \lambda_{-t}(L)^{-1}
\]
\[
  \lambda_t(L) = 1 + tL
\]

\[
  \boxed{\mathcal{E}(G) \hookrightarrow R(G)} \longrightarrow R_{\uncertain{\mathbb{Q}}}(G)
\]
\note{sous $R(G)$, un $\mathbb{Z}$ biffé ; l'indice du dernier $R$ est
douteux ; de l'encadré descend un double trait vers la ligne suivante}

\[
  \sigma_t = [1 - (tE - t^{2}\lambda^{2}(E) + \ldots)]^{-1}
\]
\[
  \struck{1 + tE}\quad 1 + tE - t^{2}\lambda^{2}E, \ \ldots \ + t^{2}E^{2}
\]
\[
  1 + tE + t^{2}(E^{2} - \lambda^{2}(E)) \ \ldots
\]

\[
  \sigma^{2}(E) = E^{2} - \lambda^{2}(E)
\]
\[
  E \times E \simeq \sigma^{2}(E) + \underline{\lambda^{2}(E)}
\]
\[
  \sigma^{2}(E) = \lambda^{2}(E) + E
\]
\note{les trois lignes sont réunies à gauche par un chevron}

\[
  \boxed{2\lambda^{2}(E) + E = E^{2}}
  \qquad
  2\,\tfrac{d(d-1)}{2} + d \quad d^{2}
\]
\note{les dimensions sont écrites sous les termes de la formule encadrée}

\[
  \lambda^{i}(xy) = \ ?
\]
\[
  \mathfrak{P}_{\alpha, \beta}(E) \qquad (\alpha_i)_{i \in I}
\]
\note{ces deux formules, en colonne à droite du feuillet, la première
isolée par une accolade}

\page{19}
$G$ groupe

$k_0(G)$ anneau \struck{des} \uncertain{des} $G$-ens.\ finis virtuels

\struck{\ill{}}

\[
  k_n(G) \struck{\ill{}} = k(\underbrace{G \times \mathfrak{S}_n}_{G(n)})
\]
($n \geqslant 0$ ; on pose $\mathfrak{S}_0 = \{1\}$) on la considère comme
$k_0(G)$-module
\[
  k_*(G) \struck{\ill{}} = \bigoplus_{n \geqslant 0} k_n(G)
\]
(somme de \struck{groupes abéliens} $k_0(G)$-modules)
\note{les parenthèses $\struck{\ill{}}$ après $k_n(G)$ et $k_*(G)$ sont
biffées}

On définit dessus une structure de $k_0(G)$-algèbre (graduée), en
définissant la multiplication $*$ par
\[
  \mathrm{cl}_m(E) * \mathrm{cl}_n(F) = \mathrm{cl}_{m+n}\big((E \times F) \wedge^{\mathfrak{S}_m \times \mathfrak{S}_n} \mathfrak{S}_{m+n}\big)
\]
où $E$ est un $G(m)$ ens.\ fini, $F$ un $G(n)$ ens.\ fini, $E \times F$ est
un $G \times G \times \mathfrak{S}_m \times \mathfrak{S}_n$ ens.\ fini,
$\mathfrak{S}_m \times \mathfrak{S}_n \to \mathfrak{S}_{m+n}$ l'hom.\ bien
connu, où $(E \times F) \wedge^{\mathfrak{S}_m \times \mathfrak{S}_n}
\mathfrak{S}_{m+n}$ est un $G \times G \times \mathfrak{S}_{m+n}$ ensemble,
qu'on considère comme un $G \times \mathfrak{S}_{m+n}$ ensemble grâce à
l'application diagonale, ce qui donne un sens à la formule.

Si $E$ est un $G$-ens.\ fini, considérons
\struck{$\tau_t(E) = \sum \ldots$}
\[
  \tau^{n}(E) = \mathrm{cl}_n(\mathrm{Mon}(I_n, E)) \in k_n(G) \qquad (n \geqslant 0)
\]
où $I_n = [1, n]$, et
\[
  \tau^{*}(E) = \sum_{n \geqslant 0} \tau^{n}(E) \in 1 + \widehat{k_*(G)^{+}}
\]
\note{un indice $t$ sous l'astérisque de $\tau^{*}$ et une ou deux lettres
après $\tau^{n}(E)$ sont biffés}

(\emph{NB} $\tau^{0}(E) = 1$, $\tau^{1}(E) = \mathrm{cl}_0(E)$)

\emph{Prop} On a
\[
  \tau^{*}(E \sqcup F) = \tau^{*}(E) . \tau^{*}(F)
\]
i.e.
\[
  \tau^{n}(E \sqcup F) = \sum_{i+j=n} \tau^{i}(E)\, \tau^{j}(F)
\]

\page{20}
\emph{Dém} \struck{\ill{}}
\[
  \mathrm{Mon}(I_n, E \sqcup F) = \coprod_{H \subset I_n} \mathrm{Mon}(H, E) \times \mathrm{Mon}(\complement_{I_n} H, F)
\]
\[
  = \struck{\ill{}} \coprod_{0 \leqslant i \leqslant n} \ \underbrace{\coprod_{J \in \mathfrak{P}_i(I_n)} \mathrm{Mon}(J, E) \times \mathrm{Mon}(\complement_{I_n} J, F)}_{\text{st.\ par } \mathfrak{S}_n}
\]
Il suffit de voir que
\[
  \mathrm{cl}_n\Big(\coprod_{J \in \mathfrak{P}_i(I_n)} \mathrm{Mon}(J, E) \times \mathrm{Mon}(\complement_{I_n} J, F)\Big) = \tau^{i}(E) * \tau^{j}(F)
\]

\emph{Cor} $\exists !$
\[
  k_0(G) \xrightarrow{\ \tau^{*}\ } 1 + \widehat{k_*(G)}^{+}
\]
(hom.\ de groupes) tel que l'on ait, $\forall\, E$ $G$-ens.\ fini,
\[
  \tau^{*}(\mathrm{cl}_0(E)) = \tau^{*}(E)
\]

Soit $H \subset \mathfrak{S}_n$ (s-groupe), on définit un hom.
\[
  q_H : k_n(G) \longrightarrow k_0(G)
\]
par la condition
\[
  q_H(\mathrm{cl}_n(E)) = \mathrm{cl}_0({}^{H}\!E)
\]
\note{le $H$ est écrit en surcharge devant $E$ ; lu comme les invariants
${}^{H}\!E$, sans certitude}
et on définit une \struck{hom.} \emph{application} (polynômiale homogène de
degré $n$)
\[
  q^{n}_H : k_0(G) \longrightarrow k_0(G)
\]
par
\[
  q^{n}_H(x) = q_H\, \tau^{n}(x)
\]

\end{document}
